\documentclass[11pt]{article}
\usepackage[margin=1in]{geometry}
\usepackage{amsmath, amssymb}
\usepackage{enumitem}
\usepackage{booktabs}
\usepackage{hyperref}
\usepackage[T1]{fontenc}
\usepackage[utf8]{inputenc}

\title{Calibration-Free Stereo Depth via a Trainable Anisotropic $\gamma$-Grid\\
\large A Computer Vision Project Report}
\author{Roman Sahakyan}
\date{2026}

\begin{document}
\maketitle

\begin{abstract}
We present a stereo depth estimation pipeline that replaces the classical pinhole intrinsic
matrix with a learnable anisotropic $\gamma$-grid, trained from a handful of in-scene click
pairs whose ground-truth distances are known. The system requires no calibration target
(no checkerboard), runs in pure NumPy with hand-derived analytical gradients, and trains
in well under a second on a CPU. We benchmark it against a classical pinhole-OLS baseline
on $43$ tape-measured points spanning $1.35$–$8.75$\,m. Under leave-one-out cross-validation
the $\gamma$-grid plus a 3-parameter quadratic post-correction reduces mean absolute depth
error from $0.373$\,m (pinhole) to $0.269$\,m, a $28\%$ improvement; it cuts the maximum
single-point error from $2.56$\,m to $1.85$\,m. We document an honest overfitting gap
between in-sample and held-out evaluation and a mirror tradeoff between the two methods on
near versus far points, both of which motivate concrete future work.
\end{abstract}

% ─────────────────────────────────────────────────────────────────────────────
\section{Problem and Motivation}

The textbook recipe for stereo depth from two consumer cameras starts with intrinsic
calibration via Zhang's algorithm: roughly twenty checkerboard images per camera are
required, and the result still assumes a single radial-tangential distortion model per
camera. We pose the inverse question: \emph{given a stereo pair with unknown intrinsics,
can we recover metric depth using only a handful of in-scene click pairs whose ground-truth
distance is known?}

Our setting is operationally simple. The user supplies $\geq 15$ calibration clicks of
the form $(u_L, v_L, u_R, v_R, Z_\text{true})$, where $Z_\text{true}$ is measured with a
tape; afterwards the system triangulates any newly clicked point in real time. We
constrain the implementation to pure NumPy with no heavyweight ML framework so that the
pipeline is reproducible from a single $\sim$1000-line script and trains in well under a
second on a CPU.

This fits both the \emph{Monocular Depth \& Scale Recovery} and \emph{Structure from
Motion / 3D Reconstruction} topics from the syllabus, with the twist that absolute
scale is recovered by user-supplied known distances rather than by monocular cues.

% ─────────────────────────────────────────────────────────────────────────────
\section{Baseline: Classical Pinhole Triangulation}
\label{sec:baseline}

The classical two-view triangulator treats each pixel as a 3-D ray and solves a
skew-line intersection in least-squares form. We summarize the derivation in this
section because the same OLS solve is reused inside our $\gamma$-grid method
(\S\ref{sec:method}).

\subsection{Geometry of two camera rays}

Camera 1 is located at a 3-D point $P_1 = [x_1, y_1, z_1]^\top$ with a normalized ray
direction $\vec d_1 = [u_1, v_1, w_1]^\top$; camera 2 likewise with $P_2$ and $\vec
d_2$. A point at distance $t_i$ along ray $i$ is
\begin{align*}
L_1(t_1) &= P_1 + t_1 \vec d_1, &
L_2(t_2) &= P_2 + t_2 \vec d_2.
\end{align*}
If the lines intersect, $P_1 + t_1 \vec d_1 = P_2 + t_2 \vec d_2$, which rearranges to
\[
t_1 \vec d_1 - t_2 \vec d_2 = P_2 - P_1.
\]

\subsection{The matrix-vector equation $Ax = b$}

Expanding component-wise yields a linear system
\[
\underbrace{\begin{bmatrix} u_1 & -u_2\\ v_1 & -v_2\\ w_1 & -w_2 \end{bmatrix}}_{A}
\underbrace{\begin{bmatrix} t_1\\t_2 \end{bmatrix}}_{x}
=
\underbrace{\begin{bmatrix} x_2-x_1\\y_2-y_1\\z_2-z_1 \end{bmatrix}}_{b}.
\]
$A$ has $3$ rows and $2$ columns: overdetermined. Pixel noise guarantees the rays will
not intersect exactly, so no exact solution exists.

\subsection{L2-norm error and OLS solution}

Define the residual $e = Ax - b$ and minimize $\|e\|_2^2$:
\begin{align*}
E(x) &= (Ax-b)^\top (Ax-b) = x^\top A^\top A x - 2 b^\top A x + b^\top b.
\end{align*}
This is a convex quadratic with gradient
\[
\nabla E(x) = 2 A^\top A x - 2 A^\top b.
\]
Setting $\nabla E = 0$ yields the normal equations and the closed-form OLS solution
\begin{equation}
A^\top A x = A^\top b \quad\Longrightarrow\quad x = (A^\top A)^{-1} A^\top b.
\label{eq:ols}
\end{equation}
The optimal $(t_1, t_2)$ are obtained in one step; for educational reference we note that
the same minimum is reachable by gradient descent $x_{k+1} = x_k - \alpha (2 A^\top A x_k -
2 A^\top b)$, which converges because $E$ is convex.

\subsection{Final 3-D point and error distance}

With optimal $(t_1, t_2)$, the closest points on each ray are
\begin{align*}
P_1^\ast &= P_1 + t_1 \vec d_1, &
P_2^\ast &= P_2 + t_2 \vec d_2.
\end{align*}
Because the lines are skew, the points do not coincide; the best estimate of the true 3-D
object is the midpoint and the reconstruction error is the gap between the two:
\[
\hat P = \tfrac{1}{2}(P_1^\ast + P_2^\ast), \qquad
\epsilon = \|P_1^\ast - P_2^\ast\|_2.
\]

\subsection{Why the pinhole baseline is insufficient}

The OLS triangulator is mathematically clean but assumes (i) perfect pinhole optics with
no radial or tangential distortion, (ii) a single focal length per camera, and (iii)
\emph{known} intrinsics. For a consumer stereo rig with unknown intrinsics we can still
fit a single scalar focal length per camera plus the baseline $(d_x, d_y)$, but this
4-parameter model leaves systematic bias at the image edges where lens distortion
dominates. Our experiments (\S\ref{sec:exp}) confirm this: the pinhole baseline achieves
$0.180$\,m mean absolute error on near-image points where geometry dominates but degrades
to $0.448$\,m on far points where distortion matters more.

% ─────────────────────────────────────────────────────────────────────────────
\section{Methodology: Anisotropic $\gamma$-Grid Stereo}
\label{sec:method}

We replace the single focal length per camera with a learnable $3 \times 3$ grid of
inverse focal lengths, separately for the $x$ and $y$ image axes. The grid is bilinearly
interpolated to produce a smooth per-pixel distortion correction without any polynomial
distortion coefficients. Triangulation itself is unchanged from \S\ref{sec:baseline} ---
we modify \emph{which rays we feed in}, not \emph{how we triangulate them}.

\subsection{Pixel and normalized coordinates}

Working directly in pixel coordinates $(u, v)$ is numerically unstable for optimization.
We normalize relative to the optical center $(C_X, C_Y) = (W/2, H/2)$:
\begin{equation}
u_n = \frac{u - C_X}{C_X}, \qquad v_n = \frac{v - C_Y}{C_Y}.
\label{eq:norm}
\end{equation}
Normalized coordinates are centered at $0$ and generally fall in $[-1, 1]$, decoupling the
image resolution from the optimization space.

\subsection{Ray equation}

A pixel $(u, v)$ defines a 3-D ray shot from the camera's optical center. Our ray model is
\begin{equation}
r = \begin{bmatrix} u_n \gamma_x \\ v_n \gamma_y \\ 1 \end{bmatrix},
\label{eq:ray}
\end{equation}
where $\gamma_x$ and $\gamma_y$ are effective inverse focal lengths (in a perfect pinhole
camera $\gamma_x = C_X / f_x$). By letting $\gamma$ \emph{vary across the image} we
inherently model lens distortion without any polynomial distortion coefficients.

\subsection{Bilinear $\gamma$ interpolation}

To make $\gamma(u, v)$ smooth and cheap we tile the image with a $3 \times 3$ grid of
trainable nodes $g_{ij}$ (cell dimensions $\text{CELL\_W} = W/2$, $\text{CELL\_H} = H/2$).
Mapping $(u, v)$ to continuous grid coordinates
\begin{equation}
f_x = u / \text{CELL\_W}, \qquad f_y = v / \text{CELL\_H},
\end{equation}
we identify the top-left integer node $(i_x, i_y)$ and the fractional distances $t_x =
f_x - i_x$, $t_y = f_y - i_y$. The interpolated value is the weighted sum of the four
surrounding nodes:
\begin{equation}
g(u, v) = (1 - t_x)(1 - t_y) g_{00} + t_x (1 - t_y) g_{10} + (1 - t_x) t_y g_{01} + t_x t_y g_{11}.
\label{eq:bilinear}
\end{equation}
This guarantees $C^0$ continuity, costs four multiplications per query, and is $O(1)$.

\subsection{Anisotropic grid: $\gamma_x \neq \gamma_y$}

Traditional pinhole models assume one focal length per camera and symmetric radial
distortion. We maintain two grids per camera, one for $\gamma_x$ and one for $\gamma_y$,
yielding a parameter tensor $\Gamma$ of shape $(R, C, 2)$. The motivation is threefold:
\begin{enumerate}[leftmargin=2em]
  \item \textbf{Sensor non-uniformity:} pixels are rarely square, so $f_x \neq f_y$.
  \item \textbf{Asymmetric distortion:} manufacturing defects in lenses are not perfectly
        circular.
  \item \textbf{Decoupled gradients:} $\gamma_x$ strictly handles horizontal disparity
        (the depth signal); $\gamma_y$ strictly handles epipolar Y-alignment.
\end{enumerate}

\subsection{Stereo geometry and triangulation}

For a matched click pair $(u_L, v_L)$ in the left image and $(u_R, v_R)$ in the right,
equation~\eqref{eq:ray} produces two rays $v_L$ and $v_R$. Let $Z$ be the depth along the
left camera's axis and $Z_R$ the depth along the right; with baseline $d = [d_x, d_y, 0]^\top$
the geometric constraint $P_L = d + P_R$ becomes
\begin{equation}
Z v_L - Z_R v_R = d \quad\Longrightarrow\quad
\underbrace{\begin{bmatrix} | & | \\ v_L & -v_R \\ | & | \end{bmatrix}}_{A}
\underbrace{\begin{bmatrix} Z \\ Z_R \end{bmatrix}}_{x} = d.
\end{equation}
The overdetermined system is solved by the same OLS step from~\eqref{eq:ols}:
\begin{equation}
x = (A^\top A)^{-1} A^\top d, \qquad Z = x_1.
\label{eq:tri}
\end{equation}

\subsection{Disparity}

In our ray space, horizontal disparity is $\text{disp} = v_{Lx} - v_{Rx}$. When $Z \gg
d_x$ we have $Z \approx Z_R$ and equation~\eqref{eq:tri} gives the well-known inverse
relation
\begin{equation}
\text{disp} \approx \frac{d_x}{Z}.
\end{equation}
This explains why stereo is precise at short range (large disparity) and imprecise at long
range (disparity $\to 0$).

\subsection{Loss design}

A naive $L = (Z_\text{pred} - Z_\text{true})^2$ is unstable because $Z \propto
1/\text{disp}$: a one-pixel error at $10$\,m produces an order of magnitude larger
gradient than at $1$\,m. We therefore optimize in ray-disparity space.

\paragraph{Relative disparity loss ($L_{1D}$).} Define $\text{disp}_\text{true} = d_x /
Z_\text{true}$ and the unit-less relative error
\begin{equation}
e_x = \frac{v_{Lx} - v_{Rx}}{\text{disp}_\text{true}} - 1 = \frac{(v_{Lx} - v_{Rx}) Z_\text{true}}{d_x} - 1,
\qquad L_{1D} = \tfrac{1}{2} e_x^2.
\label{eq:l1d}
\end{equation}
This formulation treats a $5\%$ depth error at $1$\,m identically to a $5\%$ error at
$10$\,m, stabilizing optimization.

\paragraph{Epipolar Y-alignment loss ($L_Y$).} For perfect stereo the vertical ray
components should match the physical vertical baseline:
\begin{equation}
e_y = (v_{Ly} - v_{Ry}) - \frac{d_y}{Z_\text{true}}, \qquad L_Y = \tfrac{1}{2} e_y^2.
\label{eq:ly}
\end{equation}
This forces the network to learn $\gamma_y$ such that epipolar lines become horizontal in
ray space.

\paragraph{Huber robustness.} Human calibration clicks contain outliers. Standard MSE
chases them; we wrap each error in a Huber loss with $\delta = 0.05$ (a $5\%$
relative-error threshold):
\begin{equation}
L_\delta(e) =
\begin{cases}
\tfrac{1}{2} e^2,                       & |e| \leq \delta,\\
\delta \big(|e| - \tfrac{1}{2}\delta\big), & |e| > \delta,
\end{cases}
\qquad
\frac{\partial L_\delta}{\partial e} =
\begin{cases}
e,                  & |e| \leq \delta,\\
\delta \cdot \text{sign}(e), & |e| > \delta.
\end{cases}
\label{eq:huber}
\end{equation}
Beyond $\delta$ the gradient magnitude is capped at $\delta$, bounding outlier influence.

\subsection{Analytical gradients}

We do not use automatic differentiation. All gradients are derived analytically using the
chain rule. Let $D = v_{Lx} - v_{Rx}$. Inside the Huber quadratic region:
\begin{equation}
\frac{\partial L}{\partial D} = e_x \cdot \frac{Z_\text{true}}{d_x},
\qquad
\frac{\partial L}{\partial v_{Lx}} = +\frac{\partial L}{\partial D},
\qquad
\frac{\partial L}{\partial v_{Rx}} = -\frac{\partial L}{\partial D},
\end{equation}
\begin{equation}
\frac{\partial L}{\partial v_{Ly}} = e_y,
\qquad
\frac{\partial L}{\partial v_{Ry}} = -e_y.
\end{equation}
From \eqref{eq:ray}, $v_{Lx} = u_{nL} \gamma_{xL}$, so $\partial v_{Lx} / \partial \gamma_{xL} =
u_{nL}$; from \eqref{eq:bilinear}, $\gamma_{xL} = \sum_{k=1}^4 w_k g_k$ where $w_k$ are
the bilinear weights. Composing,
\begin{equation}
\frac{\partial L}{\partial g_k}
= \frac{\partial L}{\partial v_{Lx}} \cdot u_{nL} \cdot w_k.
\label{eq:dLdg}
\end{equation}
Gradients with respect to the baseline are
\begin{equation}
\frac{\partial L}{\partial d_x} = e_x \cdot \left(-\frac{D\, Z_\text{true}}{d_x^2}\right),
\qquad
\frac{\partial L}{\partial d_y} = e_y \cdot \left(-\frac{1}{Z_\text{true}}\right).
\end{equation}
Outside the Huber quadratic region, every $e_\bullet$ in the equations above is replaced
by $\delta \cdot \text{sign}(e_\bullet)$. The exact logic above is implemented in
\texttt{forward\_grads(\dots)} in \texttt{stereo\_app.py}.

\subsection{Adam optimizer}

For any parameter $\theta \in \{\Gamma_L, \Gamma_R, d\}$ with gradient $g_t$ at step $t$:
\begin{align}
m_t &= \beta_1 m_{t-1} + (1 - \beta_1) g_t, &
v_t &= \beta_2 v_{t-1} + (1 - \beta_2) g_t^2, \\
\hat m_t &= m_t / (1 - \beta_1^t), &
\hat v_t &= v_t / (1 - \beta_2^t), \\
\theta_t &= \theta_{t-1} - \eta \, \hat m_t / (\sqrt{\hat v_t} + \epsilon).
\end{align}
We use $\beta_1 = 0.9$, $\beta_2 = 0.999$, $\epsilon = 10^{-8}$. Different parameters get
different learning rates because they live on vastly different scales: $\eta_\Gamma =
10^{-1}$ (dimensionless, $\sim 0.45$) and $\eta_d = 10^{-3}$ (meters, $\sim 0.1$). A
\texttt{ReduceLROnPlateau} schedule halves all three rates after $800$ epochs without
improvement, and early stopping triggers after $3000$ stalled epochs.

\subsection{Regularization}

With $36$ grid parameters per camera ($+$ baseline) and only $\sim 15$--$40$ calibration
points, overfitting is a real risk. We add two regularization terms directly to the
gradients.

\paragraph{Anchor prior.} Pulls $\Gamma$ toward its physical initial value $G_0 = C_X /
F_\text{init}$:
\begin{equation}
\nabla_\text{anchor} = \lambda_a (\Gamma - G_0).
\end{equation}
Unused grid cells therefore stay physically sensible rather than drifting.

\paragraph{Smoothness Laplacian.} Adjacent nodes should be similar. The 4-neighbor
discrete Laplacian is
\begin{equation}
\nabla^2 G_{i,j} \approx 4 G_{i,j} - G_{i-1,j} - G_{i+1,j} - G_{i,j-1} - G_{i,j+1},
\end{equation}
and adding $\lambda_s \nabla^2 G$ to the gradient makes the grid behave like an elastic
sheet. We use $\lambda_a = \lambda_s = 10^{-3}$.

\subsection{Polynomial post-correction}

After the $\gamma$-grid converges, slight systematic biases may remain due to unmodeled
physics (lens thickness, sensor non-linearity). We absorb the residual by fitting a 1-D
quadratic on the calibration set:
\begin{equation}
Z_\text{true} = a + b\, Z_\text{pred} + c\, Z_\text{pred}^2,
\end{equation}
solved instantly by ordinary least squares on a Vandermonde matrix of $Z_\text{pred}$.
This adds $3$ parameters; on our data the fit is $(a, b, c) = (0.020, 0.944, 0.015)$.

\subsection{Numerical stability safeguards}

\begin{enumerate}[leftmargin=2em]
  \item \textbf{Parallel rays:} if $|v_{Lx} - v_{Rx}| < 10^{-9}$ the depth is undefined;
        the code returns \texttt{NaN}.
  \item \textbf{Near-singular $A^\top A$:} falls back to the 1-D disparity formula $Z =
        d_x / D$.
  \item \textbf{$\Gamma$ clipping:} values are clipped to $[0.05, 5.0]$, preventing
        negative or infinite effective focal lengths.
  \item \textbf{Minimum baseline:} $d_x \geq 0.02$\,m, preventing division by zero in
        \eqref{eq:l1d}.
\end{enumerate}

% ─────────────────────────────────────────────────────────────────────────────
\section{Auxiliary Components}

\subsection{Fundamental matrix for click assistance}

To help the user click matching points, the UI draws an epipolar line in the right image
whenever the user clicks the left. The fundamental matrix $F$ maps a left point $x_L$ to a
right line $\ell_R = F x_L$. We construct $F$ locally from the per-pixel $\gamma$ values:
\begin{equation}
K_L^{-1} = \begin{bmatrix}
\gamma_{xL}/C_X & 0 & -\gamma_{xL} \\
0 & \gamma_{yL}/C_Y & -\gamma_{yL} \\
0 & 0 & 1
\end{bmatrix},
\qquad
E = [d]_\times = \begin{bmatrix}
0 & -d_z & d_y \\
d_z & 0 & -d_x \\
-d_y & d_x & 0
\end{bmatrix},
\end{equation}
\begin{equation}
F = K_R^{-\top} E\, K_L^{-1}, \qquad
[a, b, c]^\top = F [u_L, v_L, 1]^\top,
\end{equation}
yielding the epipolar line $a u_R + b v_R + c = 0$. Because the right pixel is not yet
known we approximate $K_R^{-1}$ using the mean of $\Gamma_R$.

\subsection{Zero-mean NCC matching}

A 1-D search along the epipolar line picks the integer pixel with maximum zero-mean
normalized cross-correlation between a left patch $A$ and a right patch $B$:
\begin{equation}
\text{ZNCC}(A, B) = \frac{\sum_i (A_i - \bar A)(B_i - \bar B)}{\sqrt{\sum_i (A_i - \bar A)^2}\sqrt{\sum_i (B_i - \bar B)^2}} \in [-1, 1].
\end{equation}
Mean subtraction gives brightness invariance; variance normalization gives contrast
invariance. The peak is refined to sub-pixel precision by fitting a parabola to the best
score and its two neighbors $(s_{-1}, s_0, s_{+1})$:
\begin{equation}
\Delta = \frac{s_{-1} - s_{+1}}{2(s_{-1} - 2 s_0 + s_{+1})},
\qquad u_\text{refined} = u_\text{best} + \Delta.
\end{equation}

% ─────────────────────────────────────────────────────────────────────────────
\section{Experiments}
\label{sec:exp}

\subsection{Dataset and protocol}

We collected $43$ hand-clicked stereo pairs across six scenes with ground-truth distances
measured by tape. Z range: $1.35$--$8.75$\,m; $12$ near points (Z $< 3$\,m) and $31$ far
points (Z $\geq 3$\,m). Image resolution is $2592 \times 1944$. We compare three
configurations:
\begin{enumerate}[leftmargin=2em]
  \item[\textbf{A.}] \textbf{Pinhole OLS baseline} --- one scalar $\gamma$ per camera plus
        $(d_x, d_y)$: $4$ parameters. Trained with the same loss as our full method, so
        the comparison isolates model capacity rather than loss design.
  \item[\textbf{B.}] \textbf{$\gamma$-grid only} --- anisotropic $3{\times}3{\times}2$ grid
        per camera plus baseline: $38$ parameters.
  \item[\textbf{C.}] \textbf{$\gamma$-grid + post-correction} --- B plus the 3-parameter
        quadratic.
\end{enumerate}
We report (i) \emph{in-sample fit} --- train on all $43$, evaluate on the same $43$ --- and
(ii) \emph{leave-one-out cross-validation}: $43$ folds of train-on-$42$, test-on-$1$. LOO
is the honest generalization estimator at this dataset size. Metrics: mean absolute error
(MAE, m), mean absolute percentage error (MAPE, \%), maximum absolute error (m), plus a
near/far MAE breakdown at the $3$\,m boundary.

\subsection{Results}

\begin{table}[h]
\centering
\caption{In-sample fit (trained and evaluated on the same $43$ points).}
\label{tab:in_sample}
\begin{tabular}{lrrrrr}
\toprule
Configuration & MAE (m) & MAPE (\%) & Max (m) & Near MAE & Far MAE \\
\midrule
A) Pinhole OLS                & 0.373 & 8.39 & 2.487 & 0.181 & 0.447 \\
B) $\gamma$-grid              & 0.145 & 3.90 & 0.694 & 0.119 & 0.156 \\
C) $\gamma$-grid + post-corr. & \textbf{0.116} & \textbf{3.46} & \textbf{0.545} & \textbf{0.115} & \textbf{0.117} \\
\bottomrule
\end{tabular}
\end{table}

\begin{table}[h]
\centering
\caption{Leave-one-out cross-validation ($n = 43$).}
\label{tab:loo}
\begin{tabular}{lrrrrr}
\toprule
Configuration & MAE (m) & MAPE (\%) & Max (m) & Near MAE & Far MAE \\
\midrule
A) Pinhole OLS                & 0.373 & 8.36 & 2.563 & \textbf{0.180} & 0.448 \\
B) $\gamma$-grid              & 0.300 & 9.04 & 1.792 & 0.373 & 0.271 \\
C) $\gamma$-grid + post-corr. & \textbf{0.269} & 8.62 & 1.854 & 0.372 & \textbf{0.229} \\
\bottomrule
\end{tabular}
\end{table}

\subsection{Discussion}

\paragraph{The $\gamma$-grid captures real distortion structure.} On the in-sample fit it
reduces MAE $2.6\times$ over the pinhole baseline ($0.373 \to 0.145$\,m) and the maximum
error drops $3.6\times$ ($2.49 \to 0.69$\,m). The drop is consistent across the Z range,
indicating the per-pixel anisotropic distortion model is doing genuine work, not just
fitting noise. The saved calibration shows $\gamma_x$ values ranging from $0.28$ to
$0.54$ across the grid --- a $90\%$ spread --- confirming the grid did not collapse to a
constant. The recovered baseline $d_x = 0.233$\,m and $d_y = 0.0005$\,m matches the
physical camera mount within sub-millimeter tolerance.

\paragraph{Overfitting is also real.} The LOO MAE doubles from $0.145$\,m (in-sample) to
$0.300$\,m, the expected price of $38$ parameters fit on $\sim 42$ points per fold even
with our anchor and Laplacian regularization. The pinhole baseline is rock-stable across
the in-sample/LOO split ($0.373 \to 0.373$) because $4$ parameters cannot overfit at this
scale --- this makes the in-sample/LOO gap the cleanest evidence of $\gamma$-grid
overfitting.

\paragraph{Near versus far reveals a mirror tradeoff.} The most interesting analytical
finding from Table~\ref{tab:loo} is that the two methods are opposites:
\begin{itemize}[leftmargin=2em]
  \item \textbf{Pinhole is good near, bad far} ($0.180$\,m near vs $0.448$\,m far). Near
        points have large disparities, so the geometry dominates the distortion error.
  \item \textbf{$\gamma$-grid is bad near, good far} ($0.373$\,m near vs $0.271$\,m far).
        The grid spends its capacity on far-point distortion correction where it matters
        most, and over-fits the near regime.
  \item Post-correction improves far-point MAE further ($0.271 \to 0.229$\,m) but cannot
        rescue near points --- a single global quadratic cannot undo overfitting that is
        locally heterogeneous.
\end{itemize}
This suggests two principled mitigations not pursued in this work: stronger
$\lambda_\text{anchor}$ in the near-image region, or reducing the grid to $2 \times 2$
and relying on the post-correction quadratic to capture remaining spatial variation.

% ─────────────────────────────────────────────────────────────────────────────
\section{Comparison with Existing Methods}

\begin{table}[h]
\centering
\begin{tabular}{p{3.2cm} p{3.5cm} p{3.5cm} p{3.3cm} p{1.8cm}}
\toprule
Method & Calibration input & Distortion model & Output & Library \\
\midrule
Zhang 2000 (\texttt{cv2.stereoCalibrate}) & $\sim$20 checkerboard images & $k_1, k_2, k_3, p_1, p_2$ & full intrinsics + extrinsics & OpenCV \\
OpenCV SGBM & requires prior $K, R, t$ & uses given $K$ & dense disparity map & OpenCV \\
Hartley \& Zisserman MVG, ch.\,12 & known fundamental matrix & none & sparse triangulation & --- \\
\textbf{Ours} & $\geq 15$ in-scene clicks at known Z & learned $3{\times}3{\times}2$ $\gamma$-grid & sparse depth + post-corr & pure NumPy \\
\bottomrule
\end{tabular}
\end{table}

The honest comparison: we lose the dense disparity map that SGBM produces and we require
ground-truth distances during calibration. We gain target-free deployment, an
order-of-magnitude lower parameter count than Zhang's intrinsic + distortion model, and
an anisotropic distortion model that is in principle more general than 1-D radial. Our
framework also produces an explicit fundamental matrix per click via $F = K_R^{-\top} [d]_\times K_L^{-1}$, which we use to draw the epipolar line in the right image and run
ZNCC sub-pixel matching for click-assist.

% ─────────────────────────────────────────────────────────────────────────────
\section{Failure Cases and Limitations}

\begin{enumerate}[leftmargin=2em]
  \item \textbf{Parallel rays.} When $|v_{Lx} - v_{Rx}| < 10^{-9}$ the triangulation
        system is singular and we return \texttt{NaN}. Happens for very distant or
        near-coplanar viewing geometry.
  \item \textbf{Disparity collapse at long range.} Beyond $\sim 10$\,m on our rig the
        disparity approaches the click-noise floor and depth uncertainty grows as $Z^2$
        regardless of model capacity.
  \item \textbf{Near-region overfitting.} Documented in Table~\ref{tab:loo} and
        \S\ref{sec:exp}. Acceptable for a click-based measurement tool but problematic for
        dense reconstruction.
  \item \textbf{Outlier clicks survive Huber.} A click on the wrong object cannot be
        rejected by any robust loss --- only by RANSAC-style consensus, which we did not
        implement.
\end{enumerate}

% ─────────────────────────────────────────────────────────────────────────────
\section{Future Work}

\begin{itemize}[leftmargin=2em]
  \item \textbf{Stronger regularization}, especially in the near-image region, to close
        the in-sample / LOO gap.
  \item \textbf{Multi-scale grids} ($3 \times 3$ refining into $9 \times 9$) for capturing
        high-frequency local lens distortions.
  \item \textbf{Uncertainty estimation} via the covariance of $x = (A^\top A)^{-1} A^\top
        d$, producing confidence intervals on $Z$.
  \item \textbf{Bundle adjustment}: jointly optimize 3-D point positions alongside camera
        parameters to minimize reprojection error.
  \item \textbf{RANSAC outlier rejection} during calibration, complementing the Huber
        loss.
\end{itemize}

% ─────────────────────────────────────────────────────────────────────────────
\section{Conclusion}

We presented a stereo depth pipeline that replaces classical pinhole intrinsics with a
trainable anisotropic $\gamma$-grid, calibrated from a handful of in-scene clicks at
known distances. The method reduces leave-one-out mean absolute depth error from
$0.373$\,m to $0.269$\,m versus a pinhole-OLS baseline on tape-measured points in
$1.35$--$8.75$\,m range, with the largest gains on far points where distortion matters
most. The same experiments expose an overfitting gap and a near/far tradeoff that
suggest concrete next steps. \medskip

\noindent\textit{Takeaway: don't calibrate the lens --- calibrate the rays.}

% ─────────────────────────────────────────────────────────────────────────────
\section*{Appendix A. Pipeline at a Glance}

\begin{enumerate}[leftmargin=2em]
  \item \textbf{Pixel click:} user selects $(u_L, v_L)$ in the left image.
  \item \textbf{Ray generation:} bilinear $\gamma$ lookup $\to$ ray $v_L$ via \eqref{eq:ray}.
  \item \textbf{Epipolar search:} build local fundamental matrix $F$; ZNCC along the line
        in the right image $\to (u_R, v_R)$ with sub-pixel refinement.
  \item \textbf{Triangulation:} OLS solve \eqref{eq:tri} $\to$ raw $Z$.
  \item \textbf{Correction:} quadratic $Z_\text{true} = a + b Z + c Z^2$ $\to$ final
        physical distance.
\end{enumerate}

\section*{Appendix B. Computational Complexity}

\begin{itemize}[leftmargin=2em]
  \item Bilinear interpolation: $O(1)$ per point, 4 multiplications.
  \item Triangulation: $O(1)$ via a $2 \times 2$ matrix solve.
  \item Training: $O(N \cdot E)$ for $N$ points and $E$ epochs. With $N \approx 40$ and
        $E \leq 20000$, total operations are $\sim 10^6$, sub-second on a modern CPU.
  \item ZNCC matching: $O(S \cdot P^2)$ where $S \approx 300$ px is the search range and
        $P = 15$ the patch size $\approx 67{,}000$ ops per click.
\end{itemize}

\section*{Appendix C. Function $\leftrightarrow$ Math Mapping}

\begin{itemize}[leftmargin=2em]
  \item \texttt{bilinear(u, v, Gamma)} $\leftrightarrow$ Equation~\eqref{eq:bilinear}.
  \item \texttt{build\_ray(u, v, Gamma)} $\leftrightarrow$ Equations~\eqref{eq:norm} and \eqref{eq:ray}.
  \item \texttt{triangulate(\dots)} $\leftrightarrow$ Equation~\eqref{eq:tri}: builds $A$ and calls \texttt{np.linalg.solve(AtA, Atd)}.
  \item \texttt{forward\_grads(\dots)} $\leftrightarrow$ Equations~\eqref{eq:l1d}--\eqref{eq:dLdg}: relative-disparity loss, Huber clipping, chain rule to $\Gamma_L, \Gamma_R, d$.
  \item \texttt{train(\dots)} $\leftrightarrow$ Adam optimizer and regularization.
  \item \texttt{fundamental\_matrix\_local(\dots)} $\leftrightarrow$ $K_L^{-1}$, $E$, $F$.
  \item \texttt{epipolar\_guided\_match(\dots)} $\leftrightarrow$ ZNCC and parabolic refinement.
  \item \texttt{fit\_postcorrection(\dots)} $\leftrightarrow$ Vandermonde solve for $(a, b, c)$.
\end{itemize}

\section*{Reproducibility}

\begin{itemize}[leftmargin=2em]
  \item Code: \texttt{stereo\_app.py} (full system), \texttt{experiments/run\_experiments.py} (this study).
  \item Data: \texttt{calibration\_points.json} ($43$ clicks), \texttt{stereo\_calibration.json} (saved $\gamma$-grid).
  \item Math reference: \texttt{docs/baseline\_classical\_stereo.tex}, \texttt{docs/final\_gamma\_grid\_method.tex}.
  \item Raw numerical results: \texttt{experiments/results.json}.
  \item Seed: \texttt{0} (numpy default RNG and \texttt{VAL\_SEED = 0}); LOO is deterministic.
\end{itemize}

\end{document}
